\documentclass{beamer}
\input{tbook-deck}
\title{Integration and divisibility: the same sign}
\author{Haobo Ma\inst{1} \and Wenlin Zhang\inst{1,2}}
\institute{\inst{1} ChronoAI Pte Ltd, Singapore\and\inst{2} National University of Singapore, Singapore}
\date{Version 1}
\begin{document}
% Manuscript: Section 1; title and attribution
\begin{frame}
\titlepage
\note{Begin with a hand-computable question. The paper authors are Haobo Ma and Wenlin Zhang; Lang introduced the sequences and Browne proposed the formula. This talk explains version 1 and its mechanism.}
\end{frame}
% Manuscript: Sections 1–3
\begin{frame}{Route}
\tableofcontents
\note{Build the two constructions, turn the floor product into a count, and compare parity. The reading edition retains the complete proofs.}
\end{frame}
\section{Question}
% Manuscript: Section 1, Theorem 1
\begin{frame}{The same answer}
For every positive integer $n$:
\begin{itemize}
\item Integrate $T_{n+1}$ and keep its reduced signed numerator.
\item Flip a sign for each $k\in\{1,\ldots,n\}$ with $n\mid k^n$.
\end{itemize}
\begin{block}{Answer}
Both are $+1$ when $4\mid n$, and $-1$ otherwise.
\end{block}
\note{This is the plain-language conclusion. Denominators are positive and numerators retain their signs. We will state the two exact equalities shortly.}
\end{frame}
% Manuscript: Section 1, defining recurrence; Proposition 2
\begin{frame}{Calculate $n=4$}
\[
T_0=1,\quad T_1=x,\qquad T_{m+1}=2xT_m-T_{m-1}\ (m\ge1).
\]
\[
T_5=16x^5-20x^3+5x,\qquad I_5=\frac16.
\]
\begin{itemize}
\item The numerator is $+1$.
\item The successes for $4\mid k^4$ are $2,4$.
\item Two negative factors: $(-1)^2=+1$.
\end{itemize}
\note{The integral is $16/6-20/4+5/2=1/6$. Emphasize the degree $n+1$. One example is not a proof; next explain the general mechanism.}
\end{frame}
\section{Counting}
% Manuscript: Proposition 3, proof, indicator identity
\begin{frame}{Floors detect divisibility}
For integer $t$ and positive integer $n$, set
\[
\delta_n(t)=\left\lfloor\frac tn\right\rfloor-\left\lfloor\frac{t-1}{n}\right\rfloor.
\]
\[
\delta_n(t)=\begin{cases}1,&n\mid t,\\0,&n\nmid t.\end{cases}
\]
Write $t=qn+u$, $0\le u<n$:
\begin{itemize}
\item At $u=0$, the floors are $q,q-1$.
\item At $u>0$, both floors are $q$.
\end{itemize}
\note{All divisions inside floors are rational. Subtract $1$ before the second division. Euclidean division also applies to negative $t$; at zero the floors are $0,-1$.}
\end{frame}
% Manuscript: Section 3; proof of Theorem 1
\begin{frame}{From radicals to a count}
The radical is the product of distinct prime factors; $\operatorname{rad}(1)=1$.
\[
n\mid k^n\quad\Longleftrightarrow\quad\operatorname{rad}(n)\mid k.
\]
\[
P_n(n;0)=(-1)^{n/\operatorname{rad}(n)}.
\]
\[
P_r(n;s)=\prod_{k=s+1}^{s+n}(1-2\delta_n(k^r)).
\]
For $12=2^2\cdot3$, the radical is $6$ and the successes are $6,12$.
\note{If $p^e\mid n$, then $e\le n$, so one copy of $p$ in $k$ supplies enough copies in $k^n$. A missing prime cannot be created. Negative factors correspond exactly to multiples of the radical.}
\end{frame}
% Manuscript: Theorem 1, public two-equality assertion
\begin{frame}{Browne's identity}
Let $a(m)$ be the reduced signed numerator of $I_m=\int_0^1T_m(x)\,dx$.
\begin{bookresult}{thm-browne}
For every positive integer $n$,
\[
\begin{aligned}
a(n+1)&=P_n(n;0),\\
P_n(n;0)&=(-1)^{n/\operatorname{rad}(n)}.
\end{aligned}
\]
\end{bookresult}
\note{The two equalities in this box match the public formal statement. Lean represents the integral by its rational coefficient sum, agreeing by termwise integration. At $n=1$ the product is $-1$ and $I_2=-1/3$.}
\end{frame}
\section{Integration}
% Manuscript: Proposition 2 and proof
\begin{frame}{Why integration gives a four-step pattern}
\[
T_m(\cos\theta)=\cos(m\theta),\qquad
I_m=\int_0^{\pi/2}\cos(m\theta)\sin\theta\,d\theta.
\]
For $m\ge2$, product-to-sum gives
\[
\begin{aligned}
I_m={}&\frac{1-\cos((m+1)\pi/2)}{2(m+1)}\\
&-\frac{1-\cos((m-1)\pi/2)}{2(m-1)}.
\end{aligned}
\]
Endpoint cosines cycle through $1,0,-1,0$.
\note{The substitution $x=\cos\theta$ reverses the endpoints, cancelling the minus sign in $dx$. Do not substitute $m=1$ into this intermediate formula: $m-1=0$. Integrate it directly.}
\end{frame}
% Manuscript: Proposition 2
\begin{frame}{The complete rational formulas}
\small
\begin{bookresult}{prop-integrals}
$I_0=1$, $I_1=1/2$; for $m\ge2$,
\[
I_m=\begin{cases}-1/(m^2-1),&m\text{ even},\\1/(m+1),&m\equiv1\pmod4,\\-1/(m-1),&m\equiv3\pmod4.\end{cases}
\]
$a(0)=1$; for $m\ge1$, $a(m)=1$ iff $m\equiv1\pmod4$, and is $-1$ otherwise.
$d(0)=1$; the positive denominators are respectively $m^2-1,m+1,m-1$, with $m=1$ included in the odd branch.
\end{bookresult}
\note{Numerators are $\pm1$ and denominators are positive, so the fractions are reduced. Thus $a(n+1)=1$ iff $4\mid n$. Reuse the existing positive-index row-sum proof for the remaining complete initial-value and denominator interface.}
\end{frame}
\section{Extensions}
% Manuscript: Proposition 3 and discussion n=16
\begin{frame}{Different counts, the same parity}
\[
n=16=2^4,\qquad D_r(16)=2^{\lceil4/r\rceil}.
\]
\begin{center}
\begin{tabular}{c|c|c|c}
$r$&$D_r(16)$&$C_r(16)$&$P_r(16;0)$\\\hline
$2$&$4$&$4$&$+1$\\
$4$&$2$&$8$&$+1$
\end{tabular}
\end{center}
$C_r$ counts successes; $P_r$ records only their parity.
\note{For squares, the successes are $4,8,12,16$; for fourth powers, all even integers succeed. Equal products do not imply equal counts.}
\end{frame}
% Manuscript: Proposition 3, complete statement
\begin{frame}{Arbitrary exponents and integer shifts}
\small
\begin{bookresult}{prop-roots}
For $n,r\ge1$, write $n=\prod_{p\mid n}p^{e_p}$ and set
\[
D_r(n)=\prod_{p\mid n}p^{\lceil e_p/r\rceil},\qquad C_r(n)=n/D_r(n).
\]
Empty products are $1$. For every integer $s$,
\[
\begin{aligned}
\sum_{k=s+1}^{s+n}\delta_n(k^r)&=C_r(n),\\
P_r(n;s)&=(-1)^{C_r(n)}.
\end{aligned}
\]
For $n>1$, $C_r(n)=n/\operatorname{rad}(n)$ iff $r\ge\max_{p\mid n}e_p$; for $n=1$ this holds for all $r\ge1$. For every $r\ge2$, $P_r(n;s)=a(n+1)$.
\end{bookresult}
\note{The two formulas give the count and product separately. The largest prime exponent controls saturation of the count, whereas $r\ge2$ already fixes the sign. Negative integers and zero belong to the statement.}
\end{frame}
% Manuscript: Proposition 3, proof
\begin{frame}{Two decisive proof steps}
\[
n\mid k^r\ \Longleftrightarrow\ D_r(n)\mid k.
\]
Since $D_r(n)\mid n$, every block of $n$ integers has $n/D_r(n)$ multiples.
\[
v_2(C_r(n))=t-\lceil t/r\rceil\quad(t=v_2(n)\ge1).
\]
For $r\ge2$: $t=1$ gives zero; $t\ge2$ gives a positive value.
\begin{alertblock}{Boundary}
At $r=1$ the product is always $-1$; $n=4$ shows why the restriction matters.
\end{alertblock}
\note{For nonzero $k$, each prime exponent satisfies $r v_p(k)\ge e_p$; check zero directly. For odd $n$ the count is odd. For $t\ge2$, use $\lceil t/r\rceil\le\lceil t/2\rceil\le t-1$.}
\end{frame}
\section{Sources}
% Manuscript: Section 4, coefficient corrections
\begin{frame}{Two coefficient corrections}
For coefficients of the zero-constant primitive,
\[
b_{m,j}=c_{m,j-1}/j\qquad(1\le j\le m+1).
\]
\begin{itemize}
\item The A164658 column multiplier should be $j/2^{j-1}$.
\item $T_7=64x^7-112x^5+56x^3-7x$, so $b_{7,6}=-56/3$, not $56/3$.
\end{itemize}
No new closed formulas are given for the separate sums A164662 and A164663.
\note{The original table displays column $m=5$ for what is denoted here by $b_{7,6}$. The printed column multiplier $j/2^{j-2}$ is twice the required value. Separate numerator sums differ from taking the numerator after summing rational numbers.}
\end{frame}
% Manuscript: Sections 1, 3 and bibliography
\begin{frame}{Sources and contributions}
\small
\begin{itemize}
\item W. Lang introduced the sequences in 2009; A. Browne proposed the formula in 2016.
\item An earlier complete Lean proof establishes the original identity.
\item This paper: self-contained explanation, exponent and interval extensions, coefficient corrections.
\item S. Finch, \emph{Idempotents and Nilpotents Modulo n}, \href{https://arxiv.org/abs/math/0605019v2}{Section 2}; NIST DLMF \href{https://dlmf.nist.gov/18.5.E1}{18.5.1}.
\item A000188--A000190 give prior root counts; S. Fried (2025), Theorem 1, proves the related Yanev identity.
\end{itemize}
\href{https://doi.org/10.5281/zenodo.23884883}{Haobo Ma and Wenlin Zhang, version 1, CC BY 4.0}.
\note{Distinguish the original question, earlier formal proof and this paper\textquotesingle s written extensions. Retain cited authors\textquotesingle{} recorded initials.}
\end{frame}
% Manuscript: Section 5, formal scope; Propositions 2–3
\begin{frame}{Reusable next steps}
\begin{itemize}
\item Original two equalities: an existing public Lean statement.
\item \href{https://github.com/the-omega-institute/trureturing-requests/issues/234}{Complete initial values and denominators}: reuse the row-sum evaluation.
\item \href{https://github.com/the-omega-institute/trureturing-requests/issues/235}{Arbitrary powers and shifted counts}: extend the divisibility construction.
\end{itemize}
The shared mechanism: divisibility switch $\longrightarrow$ count $\longrightarrow$ parity.
\note{Both tasks have written proofs. The counting extension is independent of integral evaluation; only the final numerator identification uses its sign. Finite computation does not replace all-index proof.}
\end{frame}
\appendix
% Manuscript: Section 5, formal and written evidence
\FormalVerification
\end{document}
